\scanpage{276}
\mspage{192}
% archive sheet number: 275

\textcircled{b} Pour tt $X\in\mathrm{Ob}\,M$, la flèche can. 4.8.2.2.\inserted{est iso,}
\inserted{i.e.} fait de $X$ une $\varinjlim$ du foncteur
canonique $C_{/\unsure{X|x}}\longrightarrow M$.

\textcircled{c} Le foncteur $\varphi$ (4.8.2.1) est pleinement
fidèle.

On a alors les implications
\diagram{\begin{tikzcd}[ampersand replacement=\&]
 b \arrow[r,Leftrightarrow] \& c \arrow[r,Rightarrow] \& a,
\end{tikzcd}}
et \deleted{\unsure{si|s}}si on suppose $i$ pl. fidèle, alors les
trois conditions $a,b,c$ sont équivalentes.

[\uline{NB} On dit que le foncteur $i:C\longrightarrow M$
est \uline{strictement générateur}, si les conditions
équivalentes $b,c$ sont satisfaites. Si $C$ est
une sous-catégorie de $M$ (\unsure{pas|non} néc. pleine),
on dit que $C$ est \uline{strictement génératrice}
si le foncteur d'inclusion est strictement
générateur, i.e. satisfait aux conditions
équivalentes $b)$, $c)$. Pour \deleted{$C$ est une \unsure{sous-c|sous-c.}}
\inserted{\unsure{ce,|en,|oui,} il faut donc}
que la condition $a)$ soit satisfaite i.e.
que pour tt $X$ dans $M$ la famille
des flèches $x\longrightarrow X$ de source dans
$C$, de but $X$, soit épimorphique stricte,
et cette condition est aussi suffisante
quand $C$ est sous-catégorie pleine de $M$.]

\margin{descend après $2^{\mathrm o}$}
$3^{\mathrm o})$ Considérons un diagramme commutatif
de foncteurs de $U$-catégories\,:
